% =============================================================================
%  05_cones.tex —— §5 锥与三角
% =============================================================================
\section{锥与三角}\label{sec:cone}

在拓扑空间（或单纯复形、胞腔复形）的同伦范畴里，没有核与余核的概念。事实上，
标准的柱构造表明，任何映射 $f\colon X\to Y$ 都同伦于一个包含
$$
X\longrightarrow \cyl(f)=Y\sqcup(X\times[0,1])\big/\!\sim,\qquad
f(x)\sim(x,1),
$$
而路径空间构造又表明它也同伦于一个纤维化。

对某些\emph{固定}的映射 $f\colon X\to Y$（而不是同伦映射类），我们可以谈论核
（若 $f$ 是纤维化，就是它的纤维）或余核（$Y/X$，即把 $X$ 的像收缩成一点所得的
空间，只要 $f$ 是上纤维化——在我们的讨论中就是包含）。一般情形下两者都不成立，
但有一样东西同时充当二者，那就是锥与 Dold--Puppe 构造。

映射 $f\colon X\to Y$ 上的\term{锥}{cone} $C_f$，是从 $Y\sqcup(X\times[0,1])$
出发，把 $X\times\{1\}$ 与其像 $f(X)\subset Y$ 等同、再把 $X\times\{0\}$ 收缩成
一点所得的空间。它嵌入映射序列
\begin{figure}[htbp]
\centering
\begin{tikzpicture}[scale=1,>=Stealth]
  % X
  \fill[gray!55] (0,0) circle (0.30);
  \draw (0,0) circle (0.30);
  % Y
  \draw (2.5,0) ellipse (0.46 and 0.78);
  % C_f
  \draw (5.4,0) ellipse (0.62 and 0.92);
  \fill[gray!55] (5.4,0) circle (0.30);
  \draw (5.4,0) circle (0.30);
  \draw (4.40,0) -- (4.93,0.63);
  \draw (4.40,0) -- (4.93,-0.63);
  % arrows
  \draw[->] (0.36,0) -- (1.98,0) node[midway,above]{\small $f$};
  \draw[->] (3.02,0) -- (4.45,0);
  % labels
  \node at (0,-1.05) {\small $X$};
  \node at (2.5,-1.05) {\small $Y$};
  \node at (5.4,-1.30) {\small $C_f$};
\end{tikzpicture}
\caption{锥 $C_f$：把 $X$ 的柱体粘到 $Y$ 上，再把 $X\times\{0\}$ 捏成一点。
灰色实心块代表 $X$ 及其在 $C_f$ 中的像（原图：\texttt{cone.eps}）}
\label{fig:cone}
\end{figure}

显然它可以充当余核：若 $X\xrightarrow{f}Y$ 是包含，则上图中的 $C_f$ 显然同伦
等价于 $Y/X$。事实上现在可以迭代这个过程：对自然的包含 $i\colon Y\to C_f$ 再取
锥，便得到序列
\begin{figure}[htbp]
\centering
\begin{tikzpicture}[scale=1,>=Stealth]
  \fill[gray!55] (0,0) circle (0.26);
  \draw (0,0) circle (0.26);
  \draw (2.0,0) ellipse (0.40 and 0.66);
  \draw (4.3,0) ellipse (0.54 and 0.80);
  \fill[gray!55] (4.3,0) circle (0.26);
  \draw (4.3,0) circle (0.26);
  \draw (3.44,0) -- (3.90,0.54);
  \draw (3.44,0) -- (3.90,-0.54);
  \draw (6.5,0) ellipse (0.58 and 0.84);
  \fill[gray!55] (6.5,0) circle (0.26);
  \draw (6.5,0) circle (0.26);
  \draw (5.58,0) -- (6.06,0.58);
  \draw (5.58,0) -- (6.06,-0.58);
  \draw (7.26,0) -- (7.74,0.60);
  \draw (7.26,0) -- (7.74,-0.60);
  \draw[->] (0.30,0) -- (1.53,0) node[midway,above]{\small $f$};
  \draw[->] (2.45,0) -- (3.47,0) node[midway,above]{\small $i$};
  \draw[->] (4.92,0) -- (5.60,0);
  \node at (0,-0.95) {\small $X$};
  \node at (2.0,-0.95) {\small $Y$};
  \node at (4.3,-1.12) {\small $C_f$};
  \node at (6.5,-1.20) {\small $C_i$};
\end{tikzpicture}
\caption{Dold--Puppe 序列：对 $i\colon Y\to C_f$ 再取锥，得到 $C_i$；它同伦于
$X$ 的悬挂 $\Sigma X$（原图：\texttt{doldpup.eps}）}
\label{fig:doldpuppe}
\end{figure}

把上图中最后一项里的 $Y$ 收缩到 $C_i$ 的右顶点，便看出 $C_i$ 同伦于 $\Sigma X$，
即 $X$ 的\term{悬挂}{suspension}。于是同伦意义下得到序列
\begin{equation}\label{eq:sp}
X\to Y\to Y/X\to\Sigma X\to\ldots
\end{equation}

对每一项取第 $i$ 个同调 $H_i$，并利用悬挂同构
$H_i(\Sigma X)\cong H_{i-1}(X)$，得到序列
\begin{equation}\label{eq:les}
H_i(X)\to H_i(Y)\to H_i(Y,X)\to H_{i-1}(X)\to H_{i-1}(Y)\to\ldots
\end{equation}
它正是与对子 $X\subset Y$ 相伴的长正合列。（一个有趣的练习：验证上面诱导的映射
$H_i(Y/X)\to H_i(\Sigma X)$ 确实是边界映射
$H_i(Y,X)\xrightarrow{\partial_*}H_{i-1}(X)$。）

在同伦意义下可以把 \eqref{eq:sp} 变成一列单纯映射；于是取相应的链复形，就得到
同调长正合列 \eqref{eq:les} 在复形层面的提升。当然，这对一切映射 $f$ 都成立，
不只对包含，只要把 $Y/X$ 换成 $C_f$。

例如若 $f$ 是纤维化，则 $C_f$ 充当映射的“核”即纤维。最极端的例子是
$f\colon X\to$ 点，此时 $C_f=\Sigma X$，即纤维 $X$ 的悬挂；在同调上，这就是把
$X$ 的次数整体平移 $1$。

所以 $C_f$ 兼具余核与核的功能，因此能给出我们一直在引用的 Whitehead 定理的
简易证明：设 $f\colon X\to Y$ 是单连通空间之间诱导同调群同构的映射，则序列
$$
H_i(X)\to H_i(Y)\to H_i(C_f)\to H_{i-1}(X)\to H_{i-1}(Y)\to\ldots
$$
表明 $H_*(C_f)=0$。于是由 Hurewicz 定理，同伦群 $\pi_*(C_f)$ 也为零；再由更
著名的 Whitehead 定理，$C_f$ 同伦等价于一点。由此立刻推出 $f$ 是同伦等价。

从拓扑里最后要记下的一点是：若 $X$、$Y$ 是单纯复形，$f\colon X\to Y$ 是单纯
映射，则锥 $C_f$ 自然地成为单纯复形，其 $i$ 维单形由 $Y$ 中的 $i$ 维单形，
再加上 $X$ 中 $(i-1)$ 维单形上的锥组成。稍作思考可知，这使得相应的（为方便计，
取上同调）复形为
$$
\renewcommand{\arraystretch}{1}
C_X\udot[\,1\,]\oplus C_Y\udot
\qquad\text{带有微分}\qquad
d_{C_f}=
\begin{pmatrix}
d_X[\,1\,] & 0\\
f & d_Y
\end{pmatrix},
$$
其中 $[\,n\,]$ 表示把复形向左平移 $n$ 格。

于是在任一阿贝尔范畴 $\A$ 中，对任何链复形的映射 $f\colon A\udot\to B\udot$，
都可以用上述公式定义锥 $C_f$：把 $C_X\udot$ 换成 $A\udot$、$C_Y\udot$ 换成
$B\udot$ 即可。若 $A\udot=A$、$B\udot=B$ 是集中在 $0$ 次的链复形，则 $C_f$ 就是
复形 $\{A\xrightarrow{f}B\}$。它的第零个上同调 $h^0(C_f)=\ker f$，第一个上同调
$h^1(C_f)=\operatorname{coker} f$，于是在不同次数上兼具两者。一般情形下，它就是
$A\udot\to B\udot$ 这个二项复形的全复形。

有一个显然的映射 $i\colon B\udot\to C_f$，而 $C_i$ 正如上面的拓扑所提示的，
拟同构于 $A\udot[\,1\,]$。所以在导出范畴里得到的不是核与余核，而是“正合三角”
$$
A\udot\to B\udot\to C\udot\to A\udot[\,1\,].
$$
（证明见 \cite{GM} III 3.5——严格说来应当先修改 $B$，把它换成拟同构的
$\cyl(f)$；但一旦把拟同构可逆化，这就不打紧了。）

于是我们得到的是长正合列而非短正合列；对上式取第 $i$ 个上同调 $h^i$，就得到标准
长正合列
$$
h^i(A\udot)\to h^i(B\udot)\to h^i(C\udot)\to h^{i+1}(A\udot)\to\ldots
$$

可见 $D(\A)$ 不是阿贝尔范畴，而是\term{三角范畴}{triangulated category}的一个
例子。三角范畴是带有函子 $T$（常记作 $[\,1\,]$）和一族\term{特异三角}
{distinguished triangle}的加性范畴，这族三角满足一列公理。精确定义请参看
\cite{GM} IV 1.1；但公理要求：对范畴中一切对象 $X$，有特异三角
$$
X\xrightarrow{\id}X\to 0\to X[\,1\,],
$$
而且任何态射 $f\colon X\to Y$ 都可以补成特异三角
$$
X\to Y\to C\to X[\,1\,].
$$
此外还有五引理在导出范畴中的类比，以及一条称为\term{八面体引理}{octahedral lemma}
的三角相容性，它的 \TeX{} 化之难可以想见。

为什么又要抽象一次？因为 $D(\A)$（$\A$ 为阿贝尔范畴）所具有的这种结构，在拓扑中
本来就有（我们一路都在暗示），却不必来自任何阿贝尔范畴。也就是说，本报告通篇
描述的拓扑，会让人猜想：某个“性质好的拓扑空间”构成的范畴应当恰好具有这种结构，
而我们显然并没有把它构造为某个阿贝尔范畴的导出范畴。

事实上我们还差一点，原因是悬挂函子 $T=\Sigma$ 在空间上不可逆。若它可逆，那显然
应是环路空间函子 $\Omega$，因为由标准论证，$\Omega$ 与 $\Sigma$ 互为伴随%
\footnote{记 $X$ 到 $Y$ 的映射同伦类为 $[X,Y]$。把可逆的 $\Sigma$ 作用到
$[X,\Sigma^{-1}Y]$ 上，得到同构
$[X,\Sigma^{-1}Y]\cong[\Sigma X,Y]\cong[X,\Omega Y]$。取
$X=\Sigma^{-1}Y$ 与 $X=\Omega Y$，便得典范映射
$\Sigma^{-1}Y\rightleftarrows\Omega Y$，它们在同伦群上诱导同构（取 $X=S^n$
即见）。故 $\Sigma^{-1}Y$ 在弱同伦等价意义下就是 $\Omega Y$。}。
特别地 $\Sigma X\simeq\Omega^{-1}X$，从而可以对空间“去环路”，把它们变成无穷
环路空间。

所以毫不奇怪：\term{谱}{spectra}（本质上是无穷环路空间，再加上它们的一切环路化
与去环路化，以追踪负次数上的同伦）按照我们描述的方式构成三角范畴（尽管有大量
技术细节从略，因为涉及的空间并非单纯复形）。

从拓扑中要记下的最后一点是：把拟同构可逆化之后，同伦的映射 $X\to Y$ 会等同起来。
与同伦 $X\times[0,1]\to Y$ 复合一下，便知只需证明 $X$ 到 $X\times[0,1]$ 的两个
包含 $\iota_0$、$\iota_1$（分别对应 $X\times\{0\}$ 与 $X\times\{1\}$）被等同。
但二者都是投影 $p\colon X\times[0,1]\to X$ 的右逆：$p\iota_i=\id$。又因为它们
在同调上是同构，把拟同构可逆化之后，它们都等同于 $p^{-1}$。

代数地，在 $D(\A)$ 中可以模仿（对偶的、上同调的）证明：按单纯复形的柱上显然的
乘积胞腔结构所指示的方式，定义复形的柱。即对任何链复形 $A\udot$，考虑
$$
\renewcommand{\arraystretch}{1}
\cyl(A\udot)=A\udot\oplus A\udot[\,1\,]\oplus A\udot,
\qquad\text{带有微分}\qquad
d=
\begin{pmatrix}
d_A & 0 & 0\\
-\id & d_A[\,1\,] & \id\\
0 & 0 & d_A
\end{pmatrix}.
$$
令 $\iota_1^*$、$\iota_2^*$ 分别为到第一、第三直和项的投影，$p^*$ 为相应两个
包含之和。它们都是链映射，事实上都是拟同构，且 $\iota_i^*p^*=\id$。于是在
$D(\A)$ 中 $\iota_1^*$、$\iota_2^*$ 都等同于 $(p^*)^{-1}$，从而彼此等同。

\begin{keytakeaways}
\item 同伦范畴里没有核与余核；锥 $C_f$ 同时扮演两者的角色，并迭代出
  $X\to Y\to C_f\to C_i\simeq\Sigma X$。
\item 取同调便回到经典的长正合列；取链复形便把它提升为特异三角。
\item 三角范畴 = 加性范畴 + 平移函子 + 一族特异三角 + 公理（含八面体引理）；
  $D(\A)$ 是最重要的例子，而同伦型/谱给出“没有阿贝尔范畴”的例子。
\item 同伦映射在 $D(\A)$ 中相等：用柱 $\cyl(A\udot)=A\udot\oplus A\udot[1]\oplus A\udot$
  把两个包含 $\iota_0$、$\iota_1$ 都等同于 $p^{-1}$。
\end{keytakeaways}

\begin{yizhu}[“锥”这个名字的来历与两种锥]
拓扑的锥 $C_f$ 与链复形的锥 $C_f$ 共用一个记号，并非巧合：链复形的锥正是
$C_f=\operatorname{cone}(f)$，其公式
$d=\begin{pmatrix}d_A[1]&0\\ f&d_B\end{pmatrix}$ 就是映射柱的胞腔微分。历史上
Dold--Puppe 在 1958 年前后系统地用锥来替代核与余核，故本文称之为 Dold--Puppe
构造。注意别与代数几何里“仿射锥”的锥混淆，二者无关。
[来源：Dold--Puppe, \textit{Homologie nicht additiver Funktoren} (1958)；
Weibel 第 1.5 节；A 级]
\end{yizhu}

\begin{yizhu}[特异三角公理到底有哪几条]
Verdier 的公理（TR1--TR4）大致是：TR1 恒等映射的三角与同构封闭；TR2 旋转
不变；TR3 任何态射可补成三角；TR4 八面体公理（把两个三角“粘”成第三个）。
本文只用到“存在补全”与“旋转”，八面体引理留作黑箱。之所以说它“\TeX{} 化之难”，
是因为它要在一张图里同时画出四个三角、十二个对象。
[来源：Verdier, \textit{Des catégories dérivées des catégories abéliennes}，
Astérisque 239 (1996)，附录；A 级]
\end{yizhu}

\begin{yizhu}[为什么 $\Sigma$ 不可逆、谱从哪来]
在空间范畴中，$\Sigma X=X\wedge S^1$ 只有右伴随 $\Omega$，没有左伴随；要让它
可逆，就得形式地添入 $\Sigma^{-1}$，这正是“谱”的稳定同伦范畴的做法。谱的
同伦范畴是三角范畴，平移函子就是 $\Sigma$，且 $\Sigma$ 可逆。本文这里只是点出
“三角范畴不必来自阿贝尔范畴”，为后文 §\ref{sec:df} 的导出函子作铺垫。
[来源：Adams, \textit{Stable Homotopy and Generalised Homology}；A 级]
\end{yizhu}
